\renewcommand{\Anno}{2001}

\renewcommand{\NumeroProblema}{2}

\renewcommand{\Testo}{L'aeroporto $A$ si trova a latitudine $\theta$ e
longitudine $\alpha$; l'aeroporto $B$ si trova a latitudine $\theta$ e
longitudine $\alpha+\varphi$ (con $0<\theta<\pi/2$ e $0<\varphi<\pi$).
Un aereo parte da $A$ e fa rotta verso $B$ seguendo il cammino pi\`{u}
breve, che come \`{e} noto \`{e} il minore dei due archi di cerchio
massimo passanti per $A$ e $B$.
\begin{enumerate}
	\item Determinare, in funzione di $\theta$ e $\varphi$, la massima
	latitudine alla quale passer\`{a} l'aereo.
	
	\item Tenendo $\varphi$ fisso, discutere il caso limite $\theta\to
	0$.
	
	\item Tenendo $\theta$ fisso, discutere il caso limite $\varphi\to
	\pi$.
\end{enumerate}
	
{\small Nota.  Si approssimi la terra ad una sfera e si trascuri
l'altezza alla quale vola l'aereo.}}

\renewcommand{\Risposta}{\begin{enumerate}
		\item  La latitudine massima \`{e}
		$$\arctan\left(\frac{\tan\theta}{\cos(\varphi/2)}\right).$$
	
		\item  La latitudine massima tende a $0$.
	
		\item  La latitudine massima tende a $\pi/2$.
	\end{enumerate}}

\renewcommand{\Aiutino}{Occorre considerare 2 figure: la sezione
contenente i punti della sfera a latitudine $\theta$ e la sezione,
perpendicolare al piano equatoriale, passante per il centro della
sfera e per il punto di massima latitudine dell'orbita.  Nella prima
sezione conviene segnare i punti $A$ e $B$, il centro $O'$, il punto
medio $D$ del pi\`{u} piccolo degli archi di estremi $A$ e $B$, e
l'intersezione $C$ tra i segmenti $AB$ e $OD$.  Nella seconda sezione
conviene segnare i punti $O'$, $C$ e $D$ definiti prima, nonch\'{e} il
centro $O$.  Ora le lunghezze di $O'C$ e $OO'$ non sono difficili da
calcolare e da queste si risale immediatamente all'ampiezza
dell'angolo $D\widehat{O}O'$, che \`{e} il complementare dell'angolo
che stiamo cercando.  }

\renewcommand{\SoluzioneUfficiale}{
}

\renewcommand{\SoluzioniAlternative}{
}

\renewcommand{\Osservazioni}{
}

\renewcommand{\Commenti}{
}